% !PRINTED 537
\weberpage{537}
\begin{center}
\S.\,163.

\gesperrt{Reduction of Abelian equations to cyclic ones.}
\end{center}

We have already seen above that in a transitive \gesperrt{Abelian} group no permutation other than the identical one occurs which leaves a figure unchanged. But from this there follows yet another theorem. Let us assume that a permutation $\pi$ of such a group is decomposed into its cycles, and that the cycle which contains the fewest terms is an $r$-term one. Then $\pi^r$ will leave the terms of this cycle unchanged, and must therefore be the identical permutation. But from this it follows that all the remaining cycles of $\pi$, if there are any others, must likewise consist of $r$ terms; hence:

\gesperrt{1. A permutation of a transitive Abelian group contains only cycles with the same number of terms.}

The number $r$ of the terms of a cycle must be a divisor of $m$, if $m$ is the degree of the group, and if $m=rs$, then $s$ is the number of the $r$-term cycles of which $\pi$ consists.

If now $P$ is the group of an \gesperrt{Abelian} equation $F(x)=0$, we take any non-identical permutation $\pi$ out of $P$ and decompose it into its cycles
\weq{1}{\pi = \gamma\gamma_1\gamma_2\cdots\gamma_{s-1},}
in which each of the cycles $\gamma$ relates to $r$ roots of the equation $F(x)=0$. We shall arrange these roots in such a way that
\weq{2}{\begin{aligned}
\gamma &= (\alpha,\alpha_1,\alpha_2\ldots\alpha_{r-1})\\
\gamma_1 &= (\beta,\beta_1,\beta_2\ldots\beta_{r-1})\\
&\ \,\vdots\\
\gamma_{s-1} &= (\sigma,\sigma_1,\sigma_2\ldots\sigma_{r-1})
\end{aligned}}
holds.

If $\pi_1$ is any permutation of $P$, then on account of the permutability
\weq{3}{\pi_1^{-1}\pi\pi_1 = \pi.}

By the theorem on the formation of the permutations $\pi_1^{-1}\pi\pi_1$ (\S.\,154, 6.), $\pi$ must therefore not be changed when the permutations $\pi_1$ are carried out in the cycles of $\pi$. But since the cycles are completely determined, apart from their ar-
%CONT>>
